\chapter{What this book constructs}

A topological symmetry line in two dimensions can be deformed without
changing a correlation function, provided that it does not cross an
operator insertion. For an ordinary finite-group symmetry, every such
line has an inverse. The Haagerup example asks us to take a more general
possibility seriously: two lines can fuse into a direct sum of lines.
The simplest characteristic relation is
\[
 \rho^2=\one+\rho+\alpha\rho+\alpha^2\rho,
 \qquad \alpha^3=\one.
\]
The coefficient of the identity is one, but the product is not the
identity. There is no inverse line that undoes $\rho$. Nevertheless,
the line has well-defined junctions, acts on states, and imposes precise
constraints on local physics. Its quantum dimension is
\[
 d=\frac{3+\sqrt{13}}2,
 \qquad d^2=1+3d.
\]
Understanding how this small piece of algebra becomes a conformal field
theory is the organizing question of the book.

There are several construction problems hidden in that question. One
can construct a fusion category that describes line junctions. One can
construct a chiral algebra whose representations have a prescribed
braiding and fusion. One can combine left- and right-moving fields into
a local theory. One can construct a microscopic Hamiltonian with the
line symmetry and investigate its long-distance limit. Each is a
substantial achievement, and each requires different data. A useful
education in the subject consists partly in learning to pass between
these descriptions and partly in learning when a proposed passage has
not yet been justified.

The central example in this book begins with a positive-definite lattice
and its vertex operators. Inside the lattice vertex algebra
\[
 B=V_{E_6\oplus A_2}\subset V_{E_8}
\]
one specifies six currents, a conformal vector, and a weight-two field
$u$. Their generated vertex algebra $U$ has central charge eight.
Associativity of its products comes from the ambient lattice theory.
The difficult questions are whether this generated algebra is proper,
whether it is rational, what its representations are, and whether all
the positivity and analytic conditions needed for a unitary full CFT
hold. The research construction studied here answers these questions
under an explicit set of finite lattice identities and stated
structural and analytic inputs. We turn that argument into a sequence
of definitions, calculations, and implications that a beginning
graduate student can follow.

This is a textbook exposition of a research construction. It is not an
independent verification of every large finite lattice computation,
and it does not treat every claim in a research manuscript as an
established classification theorem. Model-specific algebraic identities
and structural facts are identified as \emph{construction inputs};
deductions from them state their assumptions. Saved numerical findings
are presented as \emph{numerical records} with their finite scope.
Many of the small examples are proved
in full on the page. This distinction makes the book usable both for
learning the mathematics and for locating where a further verification
would strengthen the construction.

\section*{The two full theories and the chain}
\addcontentsline{toc}{section}{The two full theories and the chain}

Even after the chiral algebra $U$ has been identified, there is more
than one natural full theory. The diagonal theory has Hilbert space
\[
 \HH_{\rm diag}=\bigoplus_i M_i\otimes\overline{M_i},
\]
where $M_i$ runs over the twelve simple $U$-modules. A selected
holomorphic extension gives another theory with Hilbert space
\[
 \HH_{E_8}=V_{E_8}\otimes\overline{V_{E_8}}.
\]
The selected extension includes an algebra structure and a specified
Haagerup action; an equality of unflavoured characters alone would
not choose those data. Both theories have $c_L=c_R=8$, conventionally
reported as $c=8$. The diagonal theory has scalar primaries beginning
at scaling dimension $2/3$, whereas the holomorphic extension's
nonidentity scalar primaries begin at dimension two. Thus they answer
different physical questions despite their shared central charge.

The regular Haagerup chain is a third object. Its local Hilbert space
is encoded by fusion paths, and its Hamiltonian is
\[
 H=-\sum_jP_j^\rho,
\]
with projectors built from the actual unitary $\mathcal H_3$
associator. Numerical studies motivate a long-distance theory with
central charge near two. The completed research record used in this
book does not identify that continuum theory, establish its exact
central charge, or rule out a sufficiently long crossover to a gapped
phase. We therefore teach both the evidence and the tests that remain
necessary. A correct finite matrix calculation and a controlled
continuum conclusion are different kinds of statement.

There is no contradiction between central charge eight in the chiral
construction and central charge near two in the chain investigation.
The modular representation category of the chiral algebra is the
\emph{center} of a Haagerup fusion category. Having Haagerup lines in
a full theory is a different requirement. In particular, the
vanishing chiral anomaly of a center does not constrain every full
CFT carrying those lines to have central charge divisible by eight.
This distinction will become concrete in the chapters on centers,
modular transformations, and defects.

\begin{figure}[htbp]
\centering
\begin{tikzpicture}[
 box/.style={draw=MidnightBlue!65,rounded corners=2pt,
 fill=MidnightBlue!3,text width=4.5cm,align=center,
 minimum height=12mm,font=\small},
 flow/.style={-{Stealth[length=2mm]},thick,MidnightBlue},
 note/.style={font=\scriptsize,align=center,text width=3.8cm}
]
\node[box] (lattice) at (0,0)
 {Explicit lattice fields\\inside $V_{E_6\oplus A_2}$};
\node[box] (fusion) at (7.0,0)
 {Unitary category $\mathcal H_3$\\and its associator};
\node[box] (voa) at (0,-2.4)
 {The generated VOA $U$\\proper, regular, $c=8$};
\node[box] (chain) at (7.0,-2.4)
 {Fusion-path Hamiltonian\\$H=-\sum_jP_j^\rho$};
\node[box] (center) at (0,-4.8)
 {$\Rep(U)\simeq Z(\mathcal H_3)$\\with a positive ribbon structure};
\node[box] (finite) at (7.0,-4.8)
 {Finite spectra and overlaps\\Tensor network evidence};
\node[box] (full) at (0,-7.2)
 {Diagonal full theory\\or a selected extension};
\node[box] (limit) at (7.0,-7.2)
 {Continuum identification\\A separate open problem};
\draw[flow] (lattice)--node[note,left]{Finite identities\\and structural theorems}(voa);
\draw[flow] (voa)--node[note,left]{Branching, indicators,\\and analytic positivity}(center);
\draw[flow] (center)--node[note,left]{Choose the full-theory\\algebra and sewing}(full);
\draw[flow] (fusion)--node[note,right]{Local projector\\construction}(chain);
\draw[flow] (chain)--node[note,right]{Controlled finite\\calculations}(finite);
\draw[flow,dashed] (finite)--node[note,right]{Requires a justified\\scaling limit}(limit);
\end{tikzpicture}

\caption{Two construction tracks studied in this book. The left
track is the central-charge-eight research construction under its
stated finite and structural inputs. The right track constructs
the physical chain exactly but requires further work to identify
its continuum limit. Sharing Haagerup data does not supply an
arrow equating the two full theories.}
\label{fig:reading-map}
\end{figure}

\section*{Prerequisites and a route through the book}
\addcontentsline{toc}{section}{Prerequisites and a route through the book}

We assume linear algebra, elementary group representations, complex
analysis through residues and Laurent series, and the basic language
of Hilbert spaces. Familiarity with either quantum mechanics or
abstract algebra is helpful; neither a previous CFT course nor a
previous course in tensor categories is required. Some later proofs
use results from operator algebras and vertex-algebra representation
theory that cannot be developed from first principles in a single
textbook. Before each such use, we explain the relevant objects and
state the particular implication being imported.

The first part builds a common language. Fusion spaces and changes of
association give a concrete entrance to categories. Radial quantization
and Ward identities give a concrete entrance to CFT. Lattice vertex
operators then provide a model in which many abstract formulas become
ordinary oscillator and exponential calculations. The Haagerup fusion
ring, its center, and its modular data are introduced only after these
examples have supplied a reason to care about them.

The second part follows the central-charge-eight construction. It
starts with an explicit field in an explicit lattice. It then explains
why finite low-weight calculations can support an all-weight
strong-generation argument, why positive stress identities control
representations, and how branching constrains the boundary fusion
category. The passage from vacuum positivity to complete unitarity
closes the structural argument. Characters then turn that structure
into explicit graded state spaces. Their coefficients are consequences
of the constructed algebra and its representation theory.

The third part turns the construction into observables. It derives
the spectra of the two full theories, works out normalized three- and
four-point functions, and computes defect and boundary quantities.
The chapter on sewing explains exactly which data a torus partition
function leaves undetermined. The fourth part examines cosets and
stress splittings before returning to the physical chain. It develops
finite-size scaling, overlap methods, tensor networks, and exact
gapped comparison states. These tools explain what would count as an
identification of the chain and why several attractive guesses remain
conditional or fail specified tests.

The final mathematical part returns to the construction of the
associators themselves. It explains how hyperbolic gamma functions
solve finite coefficient equations, how exact algebra and positivity
enter the noncyclic examples, and how planar diagrams reconstruct
extended Haagerup. These comparisons clarify which methods generalize
and which additional steps are still needed to obtain a CFT.

A reader mainly interested in physics can first study the CFT and
lattice chapters, the explicit construction, characters, full theories,
correlators, and the chain. The category chapters can be read alongside
the first appearance of a defect or extension. A reader mainly
interested in mathematics can first follow fusion categories, lattice
vertex algebras, finite closure, rationality, branching, and unitarity,
then use the physical chapters to see which information these
structures actually predict. Both routes benefit from doing the small
worked examples rather than treating them as illustrations.

A compact first tour can follow the Ising model across the book.
Derive its null vector in Exercise~\ref{ex:cft-null}, obtain its
three chiral weights from the affine coset in Worked
example~\ref{ex:voa-Ising-branching}, and construct its full four-point
function in Section~\ref{sec:sewing-Ising}. Worked
example~\ref{ex:defect-Ising-boundaries} supplies its fixed and free
boundaries, while Worked example~\ref{ex:scaling-Ising-calibration}
derives its continuum spectrum from a periodic spin chain.
Together these give a small complete comparison theory for the
different stages of the Haagerup problem.

The following route begins with calculations a physics reader can
carry through, then returns to the construction's existence proofs.
The chapter references also make it possible to revisit a method when
it is needed. Results imported from the construction retain the
inputs and hypotheses stated there.

\begin{center}
\small
\renewcommand{\arraystretch}{1.18}
\begin{tabularx}{\textwidth}{@{}p{.29\textwidth}X@{}}
\toprule
Read & Calculation or skill to acquire\\
\midrule
Chapter~\ref{ch:cft}; the Ising example in
Chapter~\ref{ch:scaling}
&
Compute $\|L_{-2}|0\rangle\|^2=c/2$ and extract $v=2J$,
$c=1/2$ from the exactly solved periodic Ising chain.
\\[3pt]
Chapters~\ref{ch:fusion}, \ref{ch:haagerup},
and \ref{ch:chain}
&
Turn an $F$-matrix into a local fusion projector and count
the allowed periodic path states.
\\[3pt]
Chapters~\ref{ch:vertex_algebras} and
\ref{ch:lattice_embedding}
&
Multiply two lattice exponentials, including their cocycle,
and read a current level from the double pole.
\\[3pt]
Chapters~\ref{ch:centers}--\ref{ch:modular} and
\ref{ch:finite_closure}--\ref{ch:unitarity}
&
Track what finite closure, positive stresses, branching,
and complete unitarity each establish about the constructed algebra.
\\[3pt]
Chapters~\ref{ch:characters}--\ref{ch:sewing}
&
Remove Virasoro descendants from a character, count full scalar
fields, and check the Ising four-point sewing example.
\\[3pt]
Chapter~\ref{ch:cosets}
&
Subtract commuting stresses and distinguish that calculation
from constructing the reduced state space and its full CFT.
\\[3pt]
Chapters~\ref{ch:scaling}--\ref{ch:tensor_networks}
&
Derive a stress estimate of $c$, normalize a wavefunction
overlap, and obtain an MPS truncation from Schmidt coefficients.
\\[3pt]
Chapters~\ref{ch:gapped}--\ref{ch:identification}
&
Compare a proven gapped parent with the physical Hamiltonian
and state which spectral, OPE, and symmetry data a continuum
identification still needs.
\\
\bottomrule
\end{tabularx}
\end{center}

Chapters~\ref{ch:arithmetic} and \ref{ch:related-categories}
can then be read for the exact arithmetic and reconstruction
methods behind further associator examples.

\section*{How to use the examples and exercises}
\addcontentsline{toc}{section}{How to use the examples and exercises}

The clickable list after the contents gives the title and page of
every worked example. It is useful for following one calculation
across the different parts of the book or for choosing a short
first reading.

The examples are chosen to expose mechanisms that recur in the hard
construction. Diagonalizing a fusion matrix explains quantum
dimensions and the exponential growth of a path Hilbert space.
Multiplying two lattice exponentials explains both a current bracket
and the origin of a cocycle sign. Subtracting descendants from a
character explains why a large coefficient is not a count of new
primary fields. A stress-tensor norm determines a finite-size estimator
of $c$. A small vacuum algebra shows why a symmetric Hamiltonian can
have several pure vacua. Each calculation should leave the reader
with a method that can be used elsewhere.

Exercises include short checks, derivations that complete an argument,
and problems that ask which hypothesis a tempting inference is
missing. Worked solutions to every exercise are collected in
Appendix~\ref{app:solutions}. The exercises
do not require reproducing large research computations. Where saved
numerical values are used, the point is to analyze what those values
mean, not to launch a larger diagonalization.

We use a small set of conventions consistently. Vertex operators are
expanded as
\[
 Y(a,z)=\sum_{n\in\Z}a_{(n)}z^{-n-1},
 \qquad L_n=\omega_{(n+1)}.
\]
Thus $a_{(n)}b$ has weight $\wt(a)+\wt(b)-n-1$.
For a full field, $\Delta=h+\bar h$ is its scaling dimension and
$s_{\rm spin}=h-\bar h$ its spin. We reserve $d$ for the noninvertible
Haagerup quantum dimension and often write $s=1+d^2=2+3d$ in
categorical formulas. Appendix~\ref{app:notation} records other potentially
confusing uses, including the difference between a central simple
object labelled $a$ and the invertible Haagerup object $\alpha$.

References are deliberately selective. The main exposition should be
readable without opening a paper or a program. A short source guide
records the principal literature and the repository constructions
behind the book, while a verification appendix describes the scope
of exact computations and formal proofs. Its dependency ledger in
Section~\ref{sec:verification-construction-ledger} collects the
inputs and conclusions of the central-charge-eight proof.
The purpose is to make the
mathematics understandable and its logical dependencies inspectable.
